\chapter{System Analysis}

\section{Existing System}

The conventional approach to document authentication relies on physical signatures applied to paper documents, often accompanied by notarization or witness attestation for legal weight. In partially digitized environments, organizations may scan signed paper documents into PDF format, apply simple electronic signatures that are essentially images overlaid on documents, or use password-protected files as a rudimentary access control measure. None of these approaches provide cryptographic proof of document integrity or signer identity.

Simple electronic signatures, which include scanned signature images, typed names, or checkbox confirmations, offer no protection against document tampering because the signature itself is merely a visual element that can be copied, moved, or applied to different documents. Password protection prevents unauthorized access but does not detect whether the document content has been altered by someone who possesses the password. These limitations make existing informal digital signing practices unsuitable for documents where authenticity and integrity are legally or operationally critical.

\section{Proposed System}

The proposed Document Signer System replaces these inadequate mechanisms with a cryptographically sound signing and verification pipeline accessible through a standard web browser. Upon registration, each user receives a unique RSA-2048 key pair generated by the system and stored securely in the database. When a user uploads a document and requests signing, the system computes a SHA-256 hash of the file content and signs this hash with the user's private key using RSA-PSS padding. The resulting signature, along with the document hash and metadata, is stored in the database and associated with the document record.

Verification allows any authorized user to confirm that a document's current content matches the content that was signed. The system recomputes the SHA-256 hash of the current file, verifies the stored signature against this hash using the signer's public key, and compares the computed hash with the hash recorded at signing time. Any discrepancy indicates either an invalid signature or post-signing tampering, both of which are reported clearly to the user through the web interface.

\section{Feasibility Study}

The proposed system is technically feasible using well-established open-source libraries and frameworks. The Python cryptography library provides production-grade implementations of RSA key generation, signing, and verification. The Django web framework offers mature support for user authentication, file upload handling, database ORM, and template rendering. The computational requirements are modest, as RSA-2048 operations complete in milliseconds on standard hardware, and the storage requirements scale linearly with the number of users and documents.

From an economic perspective, the system is feasible because all software components are open-source and deployable on standard server infrastructure without licensing costs. From an operational perspective, the browser-based interface eliminates the need for client-side software installation, making the system immediately usable on any device with web access. From an academic perspective, the project is well-scoped for a final-year implementation because it integrates multiple core computer science topics including cryptography, database design, web development, and software engineering principles.

\section{Comparative Analysis}

Table~\ref{tab:comparison} presents a comparison between the existing informal signing methods and the proposed cryptographic system across several evaluation criteria.

\begin{table}[H]
\centering
\caption{Comparison of Existing and Proposed Systems}
\label{tab:comparison}
\begin{tabularx}{\textwidth}{|l|X|X|}
\hline
\textbf{Criterion} & \textbf{Existing System} & \textbf{Proposed System} \\
\hline
Authenticity & Visual signature only & RSA-2048 cryptographic proof \\
\hline
Integrity Check & Not available & SHA-256 hash verification \\
\hline
Tamper Detection & Manual inspection & Automated hash comparison \\
\hline
User Interface & Paper or basic PDF tools & Web dashboard with status tracking \\
\hline
Audit Trail & Limited or absent & Database records with timestamps \\
\hline
Scalability & Poor for digital workflows & Suitable for multiple users and documents \\
\hline
\end{tabularx}
\end{table}

The comparative analysis demonstrates that the proposed system provides measurable security and usability improvements over conventional document signing practices while remaining practical for academic demonstration and small-scale organizational deployment.
